% ============================================
% SECTION 5: SUCCESS STORIES & CASE STUDIES
% ============================================
\section{Case Studies: Successes That Define the Model}

\subsection{Palantir x Airbus --- A350 Production Optimization}

The Airbus engagement represents the canonical FDE success story and
illustrates the circular flow from custom solution to platform product.

\begin{tcolorbox}[
    colback=mslightgray, colframe=msblue,
    boxrule=0.5pt, arc=3pt,
    left=8pt, right=8pt, top=6pt, bottom=6pt
]
\textbf{Context}: Palantir FDEs embedded in Airbus factories in Hamburg
and Toulouse to address A350 production bottlenecks and maintenance
inefficiencies.

\vspace{4pt}
\textbf{Results}:
\begin{itemize}[topsep=2pt]
    \item Production cycle accelerated by one-third.
    \item Maintenance cost savings of 10--20\%.
    \item 35 technical cancellations avoided.
\end{itemize}

\vspace{4pt}
\textbf{Platform impact}: Solutions built for Airbus were generalized
into the \textbf{Skywise} platform, now used by 100+ airlines globally.
This is the FDE circular flow at its most powerful: bespoke work for one
customer becomes a platform serving an entire industry.
\end{tcolorbox}

\subsection{Salesforce x Reservation Platform --- Agentforce Deployment}

A Salesforce FDE pod was deployed to a reservation platform client
experiencing issues with Agentforce integration. The FDE team served as
the bridge between the customer's operational reality and Salesforce's
product engineering team, diagnosing and resolving issues that neither
side could address alone.

\subsection{Invisible Tech x Food Delivery --- AI Classification}

FDEs from Invisible Tech, embedded with a food delivery client,
discovered critical catalog data gaps that were limiting AI agent
performance. They built an AI classification system to address the gaps,
producing a measurable increase in sales. The case illustrates how FDEs
surface problems that customers do not even know they have.

\subsection{Alumni as Founders}

The FDE pipeline produces founders at a striking rate:

\begin{itemize}
    \item \textbf{Anduril Industries}: Two co-founders are ex-Palantir
        FDEs. Now a major defense technology company.
    \item \textbf{Hex} (Barry McCardel): 5-year Palantir FDE, now CEO
        of Hex, a data collaboration platform.
    \item \textbf{Veraset}: Founded by ex-FDEs, reached \$15M ARR,
        subsequently acquired.
\end{itemize}

% ============================================
% SECTION 6: RISKS, FAILURES & CRITICISM
% ============================================
\section{Risks, Failures \& Criticism}

\subsection{Financial Scalability: The Unit Economics Problem}

The most fundamental risk of the FDE model is its economics. FDEs
represent the most expensive category of engineering talent, and their
value proposition only holds at scale:

\begin{center}
\rowcolors{2}{white}{mslightgray}
\begin{tabular}{L{4cm}L{4cm}L{4.5cm}}
\toprule
\rowcolor{msblue}
\textcolor{white}{\textbf{Contract Size}} &
\textcolor{white}{\textbf{FDE Economics}} &
\textcolor{white}{\textbf{Assessment}} \\
\midrule
\$20K--\$100K &
    ``Lighting equity on fire'' &
    \textcolor{msred}{\textbf{Destructive}} \\
\$100K--\$1M &
    Marginal, high risk &
    \textcolor{msorange}{\textbf{Caution}} \\
\$1M--\$10M &
    Viable with discipline &
    \textcolor{msyellow}{\textbf{Conditional}} \\
\$10M--\$100M+ &
    Sweet spot &
    \textcolor{msgreen}{\textbf{Optimal}} \\
\bottomrule
\end{tabular}
\end{center}

\vspace{4pt}
Ex-Palantir CFO Colin Anderson's assessment is stark: deploying FDEs on
\$20K--\$100K contracts is ``lighting equity on fire.'' The 25--40\%
compensation premium over equivalent SWEs means FDE programs must target
7+ figure contracts to generate positive ROI. Satya Nadella publicly
praised the model, yet Anderson's criticism underscores that admiration
and economics are different things.

\subsection{``Spectacular Pyres'': Palantir's Failed Pilots}

Palantir's own history includes many failed FDE deployments---internally
described as ``spectacular pyres of time, money, and travel expenses that
amounted to nothing.'' The company operated with an explicit venture
model: most FDE bets fail, but the winners pay back the losses many times
over. This requires organizational tolerance for failure that most
companies lack.

\subsection{Burnout \& Talent Churn}

The human cost is documented:

\begin{itemize}
    \item \textbf{Palantir Glassdoor}: Work-life balance rated 2.7/5
        (lowest category across all employee reviews).
    \item \textbf{Travel burden}: 3--4 days per week on-site at customer
        locations, often in different cities or countries.
    \item \textbf{Brain drain}: ``Real innovators went to OpenAI, war
        dogs to Anduril'' (Glassdoor review). The best FDEs eventually
        leave for opportunities with more autonomy or equity.
\end{itemize}

\subsection{Structural Risks}

\begin{itemize}
    \item \textbf{Bespoke sprawl}: Without platform-first discipline,
        FDE programs generate thousands of custom solutions that become
        unmaintainable.
    \item \textbf{Pilot Purgatory}: Without FDEs, customers get stuck.
        With mediocre FDEs, they get stuck more expensively.
    \item \textbf{Product flaw masking}: Over-reliance on FDE human
        intervention can mask structural product defects that should be
        fixed in the platform.
    \item \textbf{External DevEx}: Palantir's FDE model created poor
        developer experience for non-FDE users, who lacked the
        hand-holding that FDE customers received.
    \item \textbf{PD vs BD fracture}: At Palantir, Product Development
        and Business Development became culturally separate organizations
        with systemic friction.
    \item \textbf{Role ambiguity}: Outside Palantir and OpenAI, ``FDE''
        describes very different jobs---from glorified solutions
        architects to true embedded engineers.
\end{itemize}

\subsection{The Cargo-Culting Warning}

Barry McCardel's warning deserves emphasis: ``Companies are cargo-culting
FDE culture. For most: absolutely f***ing not.'' Critics argue that many
companies adopting the FDE label are simply rebranding implementation
consulting. The litmus test: if FDE-generated revenue is counted as ARR
but the customer cannot operate without ongoing FDE presence, ``you're
deluding yourself.''

\subsection{Role Degradation}

Multiple sources note that the FDE role is at risk of degradation toward
a Solutions Architect function---less prestigious, less technical, and
with declining compensation. This is particularly acute at companies that
adopt the FDE title without the operational model.
